\chapter{Đam Mê Trải Nghiệm}
\markboth{Đam Mê Trải Nghiệm}{Đam Mê Trải Nghiệm}

\begin{truyen}{Chỉ Muốn Bay Nhảy}{Experience Without Discipline}
\chuhoa{M}{ột bạn nói với vẻ rất thời thượng: ``Em không hợp kiểu học khuôn mẫu. Em muốn đi nhiều, trải nghiệm nhiều, làm thứ gì đó tự do hơn. Em học trong lớp thấy bị bó.''}

\teacherpair{Cô hỏi nhẹ: ``Đi nhiều để làm gì?''}{The teacher gently asks: ``Travel and experience for what?''}

Bạn ấy đáp rất nhanh: ``Để mở mang đầu óc ạ.''

\teacherpair{Mở mang đầu óc là tốt. Nhưng nếu em chỉ thích phần bay nhảy mà không chịu nổi phần lặp lại, rèn kỹ năng, đúng giờ, hoàn thành việc đến nơi đến chốn, thì em không đang theo đuổi trải nghiệm. Em đang theo đuổi cảm giác mới lạ. Hai thứ đó khác nhau rất xa.}{``Broadening the mind is good. But if you only like the flying-around part and cannot endure repetition, skill-building, punctuality, and finishing work properly, then you are not pursuing experience. You are pursuing novelty. Those two things are very far apart.''}

Bạn ấy ngồi yên hẳn. Vì câu này chạm đúng điểm nhiều bạn ngại nghe nhất: mình không ghét khuôn mẫu, mình ghét kỷ luật.

Cô nói thêm: ``Trải nghiệm có giá trị là thứ để lại trong em một kỹ năng, một góc nhìn, một mức chịu đựng mới. Còn nếu sau mỗi chuyến đi em chỉ mang về ảnh đẹp và thêm lý do để coi thường việc học căn bản, thì em đang tiêu thụ cuộc sống chứ chưa thật sự sống sâu với nó.''

\textit{(The student romanticizes freedom and experience. The teacher separates meaningful exposure from addiction to novelty.)}
\end{truyen}

\begin{baihoc}
Trải nghiệm chỉ nuôi lớn con người khi nó đi cùng khả năng chịu trách nhiệm, không phải chỉ cùng những bức ảnh đẹp. Tự do có giá trị nhất khi nó được chống đỡ bằng tay nghề và kỷ luật.
\textit{(Experience grows a person only when it comes with responsibility, not just beautiful photos. Freedom has the greatest value when it is supported by skill and discipline.)}
\end{baihoc}

\begin{ghinhoanh}
\textbf{Short English to remember:}

Novelty is not mastery.

Freedom still needs discipline.
\end{ghinhoanh}

\begin{tuhocnhanh}
1. Em thật sự muốn trải nghiệm, hay chỉ muốn thoát khỏi thứ đang đòi hỏi mình nghiêm túc? \textit{(Do you truly want experience, or just want to escape what demands seriousness from you?)}

2. Một kỹ năng nào em cần rèn bền hơn để tự do của mình có giá trị thật? \textit{(What skill do you need to train more steadily so your freedom gains real value?)}
\end{tuhocnhanh}

\ngancach